\label{re:review_1}

\begin{center}
  \large\bf Response to Reviewer 1
\end{center}

We would like to thank you for your valuable comments and suggestions. More importantly, we are very pleased that we can have so many beneficial discussions on multisource remote sensing image classification with you. We believe these discussions must be (and it has been) helpful for us to improve this work.

We have very carefully followed all the comments and cleared up every issue raised. We provide below a point-by-point response to the comments as follows.

\begin{tcolorbox}[title={Reviewer \#1, Comment \#1}] 
  \vspace{1em}

  \color{blue}
  The manuscript does not clearly explain how the proposed SMCB differs substantially from existing dynamic convolution or semantic-aware filtering methods. Please provide in-depth description to highlight the technical contribution of the proposed method.

  \vspace{1em}
\end{tcolorbox}

\textbf{Response:}

TODO

\begin{tcolorbox}[title={Reviewer \#1, Comment \#2}]
  \vspace{1em}
  
  \color{blue}
  Please explain why frequency-domain modulation is more robust to misalignment and provide controlled experiments with different degrees of artificial spatial shifts to verify this claim.

  \vspace{1em}
\end{tcolorbox}

\textbf{Response:}

TODO

\begin{tcolorbox}[title={Reviewer \#1, Comment \#3}]
  \vspace{1em}
  
  \color{blue}
  Eq. 11 and 12 appear to place the value features inside the Softmax operation, which is inconsistent with the standard attention formulation. Please check carefully or provide more descriptions.

  \vspace{1em}
\end{tcolorbox}

\textbf{Response:}

TODO

\begin{tcolorbox}[title={Reviewer \#1, Comment \#4}] 
  \vspace{1em}
  
  \color{blue}
  The paper states that the query and key features are multiplied to obtain an affinity map, and then a dynamic kernel is generated. However, the exact dimensions of the affinity map, the reshaping operation, and the dynamic kernel are not sufficiently clear. Please add more discussions.

  \vspace{1em}
\end{tcolorbox}

\textbf{Response:}

TODO

\begin{tcolorbox}[title={Reviewer \#1, Comment \#5}] 
  \vspace{1em}
  
  \color{blue}
  The comparison should be expanded to include more recent and stronger baselines.

  \vspace{1em}
\end{tcolorbox}

\textbf{Response:}

TODO

\begin{tcolorbox}[title={Reviewer \#1, Comment \#6}] 
  \vspace{1em}
  
  \color{blue}
  Although SGFNet achieves the best overall accuracy, it does not consistently outperform all baselines for every class. This indicates that the proposed method may still struggle with some spectrally ambiguous or minority classes. The authors should provide a deeper analysis of these failure cases.

  \vspace{1em}
\end{tcolorbox}

\textbf{Response:}

TODO


\vspace{2em}

\begin{center}
  \rule[-10pt]{14.3cm}{0.05em}
\end{center}

Thanking you again for your precious time and valuable efforts spent in reviewing our manuscript. We greatly appreciate your expertise and the dedicated efforts you have made to assist us in enhancing the quality of this manuscript. We hope that we have clarified all the issues raised by you, and our revised manuscript will again meet your expectations. We eagerly anticipate any additional insights or suggestions you might offer.

Best wishes and thanks again.

\textit{Yuwei Zhao and Feng Gao}

Ocean University of China

Email: gaofeng@ouc.edu.cn
